\documentclass[12pt, a4paper]{article}

% Packages
\usepackage[UTF8]{ctex}
\usepackage[margin=2.5cm]{geometry}
\usepackage{listings}
\usepackage{xcolor}
\usepackage{titlesec}
\usepackage{enumitem}
\usepackage{hyperref}
\usepackage{booktabs}
\usepackage{fancyhdr}
\usepackage{mdframed}
\usepackage{fontawesome5}

% Page style
\pagestyle{fancy}
\fancyhf{}
\rhead{AWS Jam Challenge}
\lhead{Amazon ECR 资源策略修复}
\cfoot{\thepage}

% Colors
\definecolor{awsorange}{RGB}{255, 153, 0}
\definecolor{codeblue}{RGB}{0, 102, 204}
\definecolor{codegray}{RGB}{245, 245, 245}
\definecolor{codecomment}{RGB}{106, 153, 85}
\definecolor{codestring}{RGB}{206, 145, 120}
\definecolor{codekeyword}{RGB}{86, 156, 214}
\definecolor{alertred}{RGB}{220, 50, 47}
\definecolor{successgreen}{RGB}{0, 128, 0}

% Section formatting
\titleformat{\section}
  {\Large\bfseries\color{awsorange}}
  {\thesection}{1em}{}[\titlerule]

\titleformat{\subsection}
  {\large\bfseries\color{codeblue}}
  {\thesubsection}{1em}{}

\titleformat{\subsubsection}
  {\normalsize\bfseries}
  {\thesubsubsection}{1em}{}

% Code listing style
\lstdefinestyle{jsonStyle}{
  backgroundcolor=\color{codegray},
  commentstyle=\color{codecomment},
  keywordstyle=\color{codekeyword}\bfseries,
  stringstyle=\color{codestring},
  basicstyle=\ttfamily\small,
  breakatwhitespace=false,
  breaklines=true,
  captionpos=b,
  keepspaces=true,
  numbers=left,
  numbersep=8pt,
  numberstyle=\tiny\color{gray},
  showspaces=false,
  showstringspaces=false,
  showtabs=false,
  tabsize=2,
  frame=single,
  rulecolor=\color{codeblue!40},
  xleftmargin=15pt,
  xrightmargin=5pt
}

\lstset{style=jsonStyle}

% Highlight box
\newmdenv[
  topline=true,
  bottomline=true,
  rightline=true,
  leftline=true,
  linecolor=awsorange,
  linewidth=1.5pt,
  backgroundcolor=awsorange!8,
  innertopmargin=8pt,
  innerbottommargin=8pt,
  innerleftmargin=10pt,
  innerrightmargin=10pt,
]{warningbox}

\newmdenv[
  linecolor=codeblue,
  linewidth=1.5pt,
  backgroundcolor=codeblue!5,
  innertopmargin=8pt,
  innerbottommargin=8pt,
  innerleftmargin=10pt,
  innerrightmargin=10pt,
]{infobox}

\newmdenv[
  linecolor=successgreen,
  linewidth=1.5pt,
  backgroundcolor=successgreen!5,
  innertopmargin=8pt,
  innerbottommargin=8pt,
  innerleftmargin=10pt,
  innerrightmargin=10pt,
]{successbox}

% Hyperlink setup
\hypersetup{
  colorlinks=true,
  linkcolor=codeblue,
  urlcolor=codeblue
}

% ============================================================
\begin{document}
% ============================================================

% Title
\begin{titlepage}
  \centering
  \vspace*{3cm}
  {\fontsize{28}{36}\bfseries\color{awsorange} AWS Jam Challenge}\par
  \vspace{0.5cm}
  {\Large\bfseries\color{black} Amazon ECR 资源策略修复}\par
  \vspace{0.5cm}
  \rule{0.6\textwidth}{1.5pt}\par
  \vspace{1cm}
  {\large 知识点总结 \& 操作指南}\par
  \vspace{2cm}
  \begin{infobox}
    \centering
    \textbf{涉及服务：} Amazon ECR \quad|\quad AWS IAM \quad|\quad AWS CodeBuild
  \end{infobox}
  \vfill
  {\small 文档生成日期：\today}
\end{titlepage}

\tableofcontents
\newpage

% ============================================================
\section{题目背景}
% ============================================================

\subsection{场景描述}

在一个典型的 CI/CD 流水线中，\textbf{AWS CodeBuild} 负责构建 Docker 镜像并将其推送至 \textbf{Amazon ECR}（弹性容器注册表）。然而，由于 ECR 仓库配置了\textbf{资源策略（Resource-based Policy）}，其中存在一条 \texttt{Deny} 规则，导致 CodeBuild 的服务角色也被拦截，无法执行推送操作。

\subsection{问题仓库}

\begin{infobox}
  \textbf{ECR 仓库名称：} \texttt{pyei-challenge-app}
\end{infobox}

\subsection{原始策略（有问题的配置）}

\begin{lstlisting}[language={}]
{
  "Version": "2012-10-17",
  "Statement": [
    {
      "Sid": "DenyPutImage",
      "Effect": "Deny",
      "Principal": "*",
      "Action": "ecr:PutImage"
    }
  ]
}
\end{lstlisting}

\begin{warningbox}
  \textbf{问题所在：} 该策略对 \texttt{Principal: "*"} 全体主体拒绝 \texttt{ecr:PutImage} 操作，
  没有任何例外，导致 CodeBuild 的服务角色（Service Role）也被拒绝推送镜像。
\end{warningbox}

% ============================================================
\section{任务目标}
% ============================================================

修改 \texttt{DenyPutImage} 语句，使用 \textbf{条件键（Condition Key）}
\texttt{aws:PrincipalArn} 排除 AWS CodeBuild 服务角色，使其不受 Deny 规则影响，
从而恢复 CI/CD 流水线的正常推送能力。

\begin{successbox}
  \textbf{验收标准：} 在 ECR 资源策略的 Deny 语句中，使用 \texttt{aws:PrincipalArn}
  全局条件键排除 CodeBuild 服务角色 ARN。
\end{successbox}

% ============================================================
\section{操作步骤}
% ============================================================

\subsection{第一步：获取 CodeBuild 服务角色 ARN}

\begin{enumerate}[leftmargin=*, label=\textbf{\arabic*.}]
  \item 登录 \textbf{AWS 管理控制台}
  \item 导航至 \textbf{IAM} $\rightarrow$ \textbf{Roles（角色）}
  \item 在搜索框输入 \texttt{codebuild} 进行筛选
  \item 点击对应的 CodeBuild 服务角色
  \item 复制其完整 ARN
\end{enumerate}

\begin{infobox}
  \textbf{ARN 格式示例：}\\
  \texttt{arn:aws:iam::<账号ID>:role/codebuild-<项目名>-service-role}
\end{infobox}

\subsection{第二步：修改 ECR 资源策略}

\begin{enumerate}[leftmargin=*, label=\textbf{\arabic*.}]
  \item 进入 \textbf{Amazon ECR 控制台}
  \item 在左侧导航栏选择 \textbf{Repositories（仓库）}
  \item 点击仓库 \texttt{pyei-challenge-app}
  \item 选择 \textbf{Permissions（权限）} 标签页
  \item 点击 \textbf{Edit policy JSON}
  \item 将策略替换为修复后的版本（见第四节）
  \item 点击 \textbf{Save（保存）}
\end{enumerate}

\subsection{第三步：验证修复结果}

在 CI/CD 流水线中重新触发构建，或通过以下 AWS CLI 命令测试推送：

\begin{lstlisting}[language=bash]
# 登录到 ECR
aws ecr get-login-password --region <region> | \
  docker login --username AWS \
  --password-stdin <account_id>.dkr.ecr.<region>.amazonaws.com

# 推送镜像
docker push <account_id>.dkr.ecr.<region>.amazonaws.com/pyei-challenge-app:latest
\end{lstlisting}

% ============================================================
\section{修复后的策略}
% ============================================================

\subsection{完整策略 JSON}

\begin{lstlisting}[language={}]
{
  "Version": "2012-10-17",
  "Statement": [
    {
      "Sid": "DenyPutImage",
      "Effect": "Deny",
      "Principal": "*",
      "Action": "ecr:PutImage",
      "Condition": {
        "ArnNotEquals": {
          "aws:PrincipalArn":
            "arn:aws:iam::123456789012:role/codebuild-xxx-service-role"
        }
      }
    }
  ]
}
\end{lstlisting}

\begin{warningbox}
  \textbf{注意：} 请将 \texttt{arn:aws:iam::123456789012:role/codebuild-xxx-service-role}
  替换为你在第一步中找到的真实 CodeBuild 服务角色 ARN。
\end{warningbox}

% ============================================================
\section{核心知识点}
% ============================================================

\subsection{IAM Policy 评估逻辑}

AWS IAM 在评估权限时遵循以下优先级规则：

\begin{center}
\begin{tabular}{clc}
  \toprule
  \textbf{优先级} & \textbf{规则} & \textbf{结果} \\
  \midrule
  1 & 显式 Deny（无条件） & 拒绝 \\
  2 & 显式 Deny（条件不满足时跳过） & 取决于条件 \\
  3 & 显式 Allow & 允许 \\
  4 & 隐式 Deny（无匹配语句） & 拒绝 \\
  \bottomrule
\end{tabular}
\end{center}

\subsection{全局条件键 \texttt{aws:PrincipalArn}}

\texttt{aws:PrincipalArn} 是一个 \textbf{IAM 全局条件键（Global Condition Key）}，用于匹配发出请求的 IAM 主体（Principal）的 ARN。

\begin{center}
\begin{tabular}{lll}
  \toprule
  \textbf{条件运算符} & \textbf{含义} & \textbf{适用场景} \\
  \midrule
  \texttt{ArnEquals} & ARN 完全匹配 & 精确排除某个角色 \\
  \texttt{ArnNotEquals} & ARN 不匹配时为真 & 排除特定角色 \\
  \texttt{ArnLike} & ARN 支持通配符匹配 & 排除一类角色 \\
  \texttt{ArnNotLike} & ARN 不匹配通配符时为真 & 排除一类角色 \\
  \bottomrule
\end{tabular}
\end{center}

\subsection{Deny + ArnNotEquals 的组合逻辑}

这是本题的核心逻辑，可用以下表格理解：

\begin{center}
\begin{tabular}{lllc}
  \toprule
  \textbf{请求者} & \textbf{ArnNotEquals 结果} & \textbf{Condition 满足？} & \textbf{Deny 生效？} \\
  \midrule
  CodeBuild 角色 & \texttt{false} & 否 & \textcolor{successgreen}{\textbf{否（允许）}} \\
  其他 IAM 主体 & \texttt{true} & 是 & \textcolor{alertred}{\textbf{是（拒绝）}} \\
  \bottomrule
\end{tabular}
\end{center}

\begin{infobox}
  \textbf{理解要点：} 条件（Condition）是 Deny 语句的\textbf{触发门槛}。只有当条件为 \texttt{true} 时，
  Deny 才生效。使用 \texttt{ArnNotEquals} 意味着：``除了指定的 ARN 以外的所有人，才触发拒绝''。
\end{infobox}

\subsection{ECR 资源策略 vs IAM 身份策略}

\begin{center}
\begin{tabular}{lll}
  \toprule
  \textbf{属性} & \textbf{ECR 资源策略} & \textbf{IAM 身份策略} \\
  \midrule
  附加对象 & ECR 仓库（资源） & IAM 用户/角色/组 \\
  控制方向 & 谁可以访问此资源 & 此主体可访问哪些资源 \\
  跨账号访问 & 支持 & 需配合资源策略 \\
  Principal 字段 & 必须指定 & 不需要（隐式为附加主体）\\
  \bottomrule
\end{tabular}
\end{center}

\subsection{Amazon ECR 常用操作权限}

\begin{center}
\begin{tabular}{ll}
  \toprule
  \textbf{权限} & \textbf{用途} \\
  \midrule
  \texttt{ecr:PutImage} & 创建或更新镜像 manifest 与 tag \\
  \texttt{ecr:GetDownloadUrlForLayer} & 拉取镜像层 \\
  \texttt{ecr:BatchGetImage} & 批量获取镜像 \\
  \texttt{ecr:GetAuthorizationToken} & 获取认证 Token（需 IAM 策略）\\
  \texttt{ecr:InitiateLayerUpload} & 初始化镜像层上传 \\
  \texttt{ecr:UploadLayerPart} & 上传镜像层分片 \\
  \texttt{ecr:CompleteLayerUpload} & 完成镜像层上传 \\
  \bottomrule
\end{tabular}
\end{center}

% ============================================================
\section{架构图解}
% ============================================================

\subsection{CI/CD 流水线与 ECR 交互流程}

\begin{infobox}
\begin{verbatim}
  开发者提交代码
       |
       v
  [AWS CodePipeline]
       |
       v
  [AWS CodeBuild]  <-- 使用 Service Role ARN
       |
       | docker push
       v
  [Amazon ECR]
  pyei-challenge-app
       |
       | ECR 资源策略检查
       v
  ┌─────────────────────────────────────────┐
  │  DenyPutImage 语句                       │
  │  Condition: ArnNotEquals                │
  │    CodeBuild Role ARN → 条件为false      │
  │    → Deny 不生效 → 允许推送 ✅           │
  └─────────────────────────────────────────┘
\end{verbatim}
\end{infobox}

% ============================================================
\section{常见错误与排查}
% ============================================================

\begin{enumerate}[leftmargin=*, label=\textbf{Q\arabic*.}]

  \item \textbf{使用了错误的条件键}\\
  \textit{错误：} 使用 \texttt{aws:PrincipalType} 或 \texttt{aws:userid}\\
  \textit{正确：} 使用 \texttt{aws:PrincipalArn}，这是按角色 ARN 做精确匹配时最直接、最清晰的全局条件键。

  \item \textbf{ARN 填写不完整}\\
  \textit{错误：} 只填写角色名 \texttt{codebuild-role}\\
  \textit{正确：} 必须填写完整 ARN，如 \texttt{arn:aws:iam::123456789:role/codebuild-role}

  \item \textbf{条件运算符选择错误}\\
  \textit{错误：} 使用 \texttt{ArnEquals}（这会导致只有 CodeBuild 被拒绝）\\
  \textit{正确：} 使用 \texttt{ArnNotEquals}（排除 CodeBuild，其余全部拒绝）

  \item \textbf{忘记保存策略}\\
  修改完 JSON 后，必须点击 \textbf{Save} 按钮，否则修改不会生效。

\end{enumerate}

% ============================================================
\section{总结}
% ============================================================

\begin{successbox}
  \textbf{核心要点回顾：}
  \begin{enumerate}
    \item ECR 资源策略的 Deny 语句可以通过 \textbf{Condition} 添加例外
    \item 使用全局条件键 \texttt{aws:PrincipalArn} 匹配请求者的 ARN
    \item \texttt{ArnNotEquals} + \texttt{Deny} 组合 = ``除了指定 ARN 以外，全部拒绝''
    \item 修复 CI/CD 流水线权限问题时，\textbf{最小权限原则}是首选方案
    \item 排查 ECR 推送失败时，需同时检查 \textbf{ECR 资源策略} 和 \textbf{IAM 身份策略}
  \end{enumerate}
\end{successbox}

\vspace{1cm}
\begin{center}
  \rule{0.4\textwidth}{0.5pt}\\
  \vspace{0.3cm}
  {\small AWS Jam Challenge 知识点总结文档}\\
  {\small 基于 Amazon ECR 资源策略排错场景}
\end{center}

\end{document}
